\documentclass[conference]{IEEEtran}

\usepackage{amsmath,amssymb}
\usepackage{bm}
\usepackage{booktabs}
\usepackage{cite}
\usepackage{graphicx}
\usepackage{siunitx}
\usepackage{url}
\usepackage[hidelinks]{hyperref}

\graphicspath{{figures/}}
\urlstyle{same}

\newcommand{\mjorbit}{MuJoCo-Orbit}
\newcommand{\mj}{MuJoCo}
\newcommand{\R}{\mathbb{R}}

\title{\mjorbit: A MuJoCo-Style Simulator for Coupled Orbital and Multibody
Spacecraft Dynamics}

\author{
\IEEEauthorblockN{John Zhang}
\IEEEauthorblockA{
Department / Lab Affiliation\\
University or Organization\\
Email: your.email@example.com}
}

\begin{document}
\maketitle

\begin{abstract}
Spacecraft operating in low Earth orbit increasingly combine orbital motion,
attitude dynamics, articulated mechanisms, contact-rich interactions, sensors,
and multiple actuator classes. Existing astrodynamics tools provide mature
orbit and environment models, while robotics engines provide efficient
rigid-body dynamics and contact simulation, but these capabilities are often
separated by incompatible state representations and frame conventions. This
paper introduces \mjorbit, a Python package that couples a MuJoCo multibody
model to a chief-orbit propagation layer through a MuJoCo-style
\texttt{MjoModel}/\texttt{MjoData} API. The simulator advances rigid and
articulated spacecraft in a local LVLH frame, propagates the chief orbit in
ECI coordinates, and assembles world-frame body wrenches from differential
gravity, J2 perturbations, drag, solar radiation pressure, magnetic torques,
gravity-gradient torques, reaction wheels, magnetorquers, control moment
gyros, and thrusters. A preliminary validation suite compares limiting cases
against Clohessy-Wiltshire relative motion, two-body and J2 orbital
propagation, torque-free rigid-body dynamics, gyrostat behavior, and coupled
articulated examples. The planned experiments will show that \mjorbit can
serve as a research platform for spacecraft robotics, proximity operations,
and attitude-control studies that need both orbital forcing and MuJoCo's
multibody/contact machinery.
\end{abstract}

\begin{IEEEkeywords}
spacecraft simulation, multibody dynamics, orbital mechanics, MuJoCo,
attitude dynamics, contact dynamics
\end{IEEEkeywords}

\section{Introduction}
\label{sec:introduction}

Low Earth orbit is becoming a setting for increasingly physical spacecraft
behaviors: robotic inspection, servicing, deployable structures, on-orbit
assembly, capture, docking, and contact with targets whose mass properties may
be uncertain. These tasks are not purely orbital, purely attitude, or purely
robotic. They require a simulator that can represent a spacecraft as an
articulated multibody system while also applying the environmental and
actuator loads that dominate orbital flight.

Astrodynamics simulators typically provide accurate orbit propagation,
environment models, and flight-dynamics workflows, but they are not designed
around general contact-rich articulated dynamics. Robotics simulators provide
fast rigid-body dynamics, constraints, actuation, and contact, but usually
assume a terrestrial inertial or gravity-fixed world. \mjorbit{} is motivated
by the gap between these two tool families. The package treats \mj{} as the
inner multibody dynamics engine and adds an orbital layer that keeps the
MuJoCo world frame aligned with a local LVLH frame attached to a propagated
chief orbit.

The current implementation exposes a compact public API:
\texttt{MjoModel} stores the compiled MuJoCo model and static orbital-coupling
metadata, \texttt{MjoData} stores one simulation rollout, and
\texttt{mjo\_forward}/\texttt{mjo\_step} mirror MuJoCo's synchronization and
stepping pattern. Static metadata is supplied through dataclasses such as
\texttt{SurfaceSpec}, \texttt{ReactionWheelSpec},
\texttt{MagnetorquerSpec}, \texttt{ControlMomentGyroSpec}, and
\texttt{ThrusterSpec}. Runtime commands and sensor measurements live on the
data object, matching the model/data split familiar to MuJoCo users.

This paper will make four main contributions. First, it describes a
MuJoCo-compatible state and stepping interface for coupled orbital and
multibody simulation. Second, it gives the wrench-coupling formulation used
to apply orbital, environmental, and actuator effects to each body in the
LVLH frame. Third, it documents implementation details that are easy to get
wrong in MuJoCo spacecraft simulations, including body-frame angular velocity,
inertia-frame orientation, and world-frame external torques. Fourth, it
presents a validation and demonstration suite covering analytical orbital
limits, attitude dynamics, actuator behavior, articulated motion, and contact.

The remainder of the paper is organized as follows. Section~\ref{sec:related}
positions \mjorbit{} relative to astrodynamics frameworks and robotics
physics engines. Section~\ref{sec:formulation} introduces the coupled state,
frames, forces, and stepping algorithm. Section~\ref{sec:experiments}
outlines the validation studies and demonstrations planned for the submission.
Section~\ref{sec:conclusion} summarizes the intended contribution and open
development items.

\section{Related Work}
\label{sec:related}

\subsection{Astrodynamics and Spacecraft Simulation}

Mission-analysis and astrodynamics tools such as GMAT, Orekit, and Basilisk
provide orbit propagation, environment models, and spacecraft dynamics
workflows for flight-dynamics analysis and control prototyping
\cite{gmat,orekit,basilisk}. These systems are natural references for
orbital fidelity, sensor modeling, and flight software integration. The
distinguishing goal of \mjorbit{} is different: it aims to make a general
MuJoCo multibody model, including articulated links and contacts, feel like a
first-class spacecraft model in orbit rather than a separate robotics
simulation wrapped around an orbital propagator.

Classical spacecraft attitude and orbital mechanics provide the analytical
limits used throughout the project. Two-body and J2 propagation, LVLH relative
motion, gravity-gradient torques, rigid-body Euler equations, gyrostat
dynamics, and actuator models are standard tools in spacecraft dynamics and
control \cite{wertz1978spacecraft,wie2008space}. \mjorbit{} does not replace
these models; instead, it embeds them as per-body and per-actuator wrench
terms applied to a MuJoCo model.

\subsection{Robotics Physics Engines}

Physics engines such as MuJoCo, Drake, and Gazebo have made articulated
rigid-body simulation, constraints, actuation, and contact widely accessible
for robotics research \cite{todorov2012mujoco,drake,gazebo}. MuJoCo is
particularly attractive for spacecraft robotics because it exposes a compact
model/data API, supports free bodies and articulated mechanisms, and provides
efficient contact dynamics. However, MuJoCo itself does not provide an orbital
reference frame, distributed environmental loads, or actuator models such as
reaction wheels, magnetorquers, and control moment gyros. \mjorbit{} keeps
MuJoCo as the dynamics backend while adding these space-specific layers.

\subsection{Positioning of This Work}

The intended contribution is therefore not a new integrator for every
astrodynamics regime, nor a replacement for high-fidelity mission analysis.
It is a simulator for research questions where the hard part is the coupling:
How does an articulated spacecraft move while the chief orbit propagates?
How do environmental torques and actuator momentum states affect the same
MuJoCo body state? How do contact events change relative motion without
leaking internal contact forces into the chief orbit? These are the questions
the formulation and experiments are designed to make concrete.

\section{Coupled Simulation Formulation}
\label{sec:formulation}

\subsection{State and Frames}

\mjorbit{} represents the chief orbit in Earth-centered inertial coordinates
as
\begin{equation}
  x_o = (\bm R,\bm V,t), \qquad \bm R,\bm V \in \R^3,
\end{equation}
with distances in kilometers and time in seconds. The MuJoCo model state
$(q,v)$ is expressed in a local LVLH frame attached to the chief. The LVLH
basis is
\begin{equation}
  \hat{\bm x}_L = \frac{\bm R}{\|\bm R\|}, \qquad
  \hat{\bm z}_L = \frac{\bm R \times \bm V}{\|\bm R \times \bm V\|}, \qquad
  \hat{\bm y}_L = \hat{\bm z}_L \times \hat{\bm x}_L ,
\end{equation}
and the rows of $C_{LI}$ map ECI vectors into LVLH coordinates. The angular
velocity of the LVLH frame is
\begin{equation}
  \bm\omega_L = C_{LI}\frac{\bm R \times \bm V}{\|\bm R\|^2}.
\end{equation}
When J2 is enabled, $\dot{\bm\omega}_L$ is computed from the current
perturbed orbital acceleration rather than assuming two-body motion.

The MuJoCo world frame is identified with the LVLH frame. For each MuJoCo
body $i$, the body center-of-mass position $\bm r_i$ and velocity
$\bm v_i$ are read from MuJoCo in SI units and converted to kilometers where
the orbital layer requires them.

\subsection{Per-Body Orbital Forcing}

The translational orbital coupling applies the exact rotating-frame apparent
acceleration to each body center of mass:
\begin{align}
  \bm a_i^L
    &= C_{LI}\left[\bm g(\bm R_i)-\bm g(\bm R)\right]
       -2\bm\omega_L \times \bm v_i^L \nonumber\\
    &\quad -\dot{\bm\omega}_L \times \bm r_i^L
       -\bm\omega_L \times (\bm\omega_L \times \bm r_i^L),
  \label{eq:apparent-acceleration}
\end{align}
where $\bm R_i = \bm R + C_{IL}\bm r_i^L$ is the absolute ECI position of
body $i$ and $\bm g(\cdot)$ includes point-mass gravity and, optionally, J2.
Equation~\eqref{eq:apparent-acceleration} is not a Clohessy-Wiltshire
linearization, though it reduces to the CW limit for small offsets in
circular two-body orbits. The force applied to body $i$ is
\begin{equation}
  \bm F_i^L = m_i \bm a_i^L .
\end{equation}

\subsection{Environmental Wrenches}

Additional external wrenches are accumulated in a per-body buffer before
being written to MuJoCo's \texttt{xfrc\_applied}. The current implementation
includes flat-plate drag and solar radiation pressure, residual magnetic
dipoles, and rotational gravity-gradient torque. For a configured surface
with body-frame normal $\hat{\bm n}$, center of pressure $\bm p$, and area
$A$, the surface module rotates the normal and center of pressure into LVLH,
computes atmosphere-relative velocity at the point, and applies force and
moment
\begin{equation}
  \bm\tau_i^L = \bm p_i^L \times \bm F_i^L.
\end{equation}
The gravity-gradient torque on a rigid body is modeled as
\begin{equation}
  \bm\tau_{gg}
    = \frac{3\mu}{\|\bm R_i\|^3}
      \hat{\bm r}_i \times \left(J_i^L \hat{\bm r}_i\right),
  \label{eq:gravity-gradient}
\end{equation}
where $J_i^L$ is the body's inertia tensor expressed in the MuJoCo world/LVLH
frame.

\subsection{Actuators and Sensors}

Actuator metadata is compiled against MuJoCo body names, while runtime
commands live on \texttt{MjoData.actuators}. Reaction wheels integrate wheel
speed from commanded torque and apply the equal-and-opposite body torque.
They also apply gyrostat coupling from stored wheel momentum. Magnetorquers
apply $\bm\tau = \bm m \times \bm B$ using the environment magnetic field.
Thrusters apply a clipped body-frame force at a body-fixed point and therefore
produce both force and torque. Single-gimbal control moment gyros are modeled
with a body-fixed gimbal axis $\hat{\bm g}$, spin axis
$\hat{\bm s}(0)$, rotor momentum $h$, and gimbal angle $\theta$:
\begin{equation}
  \hat{\bm s}(\theta)
    = \cos\theta\,\hat{\bm s}(0)
      + \sin\theta\,\left(\hat{\bm g}\times\hat{\bm s}(0)\right),
\end{equation}
with commanded output torque
\begin{equation}
  \bm\tau_{cmg}
    = -h\dot\theta\,
      \left(\hat{\bm g}\times\hat{\bm s}(\theta)\right).
\end{equation}

Sensors are resolved from the MuJoCo sensor catalog and exposed through a
runtime namespace. The package currently supports canonical MuJoCo sensor
data access, stochastic measurement helpers, and custom orbit-aware sensors
for quantities such as sun, horizon, star, and magnetic-field directions.

\subsection{Stepping Algorithm}

Each call to \texttt{mjo\_step} clears the external-wrench buffer, assembles
environmental and actuator wrenches, copies them into MuJoCo, computes the
net translational external force for chief-orbit feedback, advances the
chief orbit with RK4, refreshes frame and environment caches, and finally
steps the MuJoCo state. Internal MuJoCo forces such as joints, constraints,
and contacts are intentionally not fed back into the chief orbit; only net
external forces accumulated in the coupling layer affect the propagated
chief state.

\section{Experiments and Demonstrations}
\label{sec:experiments}

The experiments will combine unit-level validation, analytical parity
checks, and end-to-end demonstrations. The current test suite already
exercises most of these cases; the paper version should convert the strongest
ones into reproducible plots, tables, and ablations.

\subsection{Analytical Validation}

The first group of experiments will verify limiting cases where closed-form
or high-accuracy reference solutions are available.

\begin{table}[t]
\centering
\caption{Planned validation experiments for the conference submission.}
\label{tab:experiments}
\small
\begin{tabular}{@{}p{0.27\linewidth}p{0.35\linewidth}p{0.27\linewidth}@{}}
\toprule
Experiment & Purpose & Metric \\
\midrule
Two-body orbit & Verify chief RK4 propagation without perturbations & Relative
specific-energy drift over one orbit \\
J2 orbit & Verify secular nodal regression & RAAN drift error after one day \\
CW relative motion & Check LVLH per-body forcing in the small-offset circular
limit & Position and velocity error vs. CW solution \\
Torque-free rigid body & Check MuJoCo free-joint attitude conventions and
full-inertia handling & Energy, angular momentum, and SciPy Euler-equation error \\
Gyrostat and CMG & Validate internal momentum storage and reaction torques &
Angular-rate and momentum consistency \\
Gravity-gradient torque & Validate rotational environmental torque & Torque
direction and magnitude vs. analytical model \\
\bottomrule
\end{tabular}
\end{table}

For the CW test, a free body is initialized with meter-scale LVLH offsets and
centimeter-per-second relative velocities in a 400 km circular orbit. The
simulated relative trajectory will be compared against the classical
Clohessy-Wiltshire solution. This experiment is important because it tests
the full rotating-frame force computation while still admitting an analytical
reference.

For attitude dynamics, the paper will compare torque-free MuJoCo rollouts
against Euler-equation integration for diagonal and non-diagonal inertia
tensors. This case is also a useful place to document MuJoCo frame
conventions: free-joint angular velocity is stored in the body frame, while
\texttt{xfrc\_applied} torques are world-frame torques.

\subsection{Environmental and Actuator Case Studies}

The second group of experiments will demonstrate the simulator's spacecraft
specific features. Candidate studies include:

\begin{itemize}
  \item Drag and solar-radiation-pressure loads on configured flat plates,
  including eclipse transitions and projected-area effects.
  \item Magnetic residual-dipole and magnetorquer torques along a LEO orbit.
  \item Reaction-wheel spin-up, speed saturation, and gyrostat coupling.
  \item Thruster force/torque allocation and the resulting feedback into the
  propagated chief orbit.
  \item Single-gimbal CMG momentum and torque envelopes for a spacecraft
  attitude maneuver.
\end{itemize}

The current ISS analysis scripts provide a useful high-level demonstration
target. For the paper, these scripts can be turned into a more general
spacecraft case study with a bounded slew maneuver, environmental torque
accumulation, and actuator limits reported in SI units.

\subsection{Articulation, Contact, and Robotics Workloads}

The third group of demonstrations should emphasize what is difficult to show
in a traditional orbit propagator. The existing spacecraft-arm example
simulates a free-floating base with a two-link arm in orbit, checks that
internal arm motion remains finite and does not inject external force into
the chief orbit, and reports end-effector motion in the LVLH frame. A paper
figure could show the arm trajectory, base reaction motion, and chief-orbit
energy drift.

The contact experiment should show two free bodies colliding inside the
LVLH-frame MuJoCo model. The key result is separation of concerns: MuJoCo
contact impulses change the relative states of the bodies, but those internal
contact forces do not appear in the orbital external-wrench buffer. The paper
can report contact impulse magnitude, wrench-buffer magnitude, total LVLH
momentum before and after collision, and post-collision relative drift.

\subsection{Ablations and Reporting Plan}

Each demonstration should include an ablation that turns off one coupling
term and shows the expected change. Examples include disabling J2 in the
RAAN-drift case, disabling gravity-gradient torque in the attitude case, and
disabling orbit feedback for thruster maneuvers. The final paper should
report runtime, timestep, model size, and error tolerances so readers can
distinguish numerical accuracy from modeling assumptions.

\section{Conclusion and Submission Plan}
\label{sec:conclusion}

This draft positions \mjorbit{} as a bridge between astrodynamics simulation
and MuJoCo-style multibody/contact dynamics. The core argument is that many
spacecraft robotics and ADCS research problems need both sides at once:
orbital environment and actuator physics must act on the same articulated
state that produces contact, joint, and rigid-body behavior.

Before submission, the paper should replace this rough outline with
quantitative results, add figures from the validation suite, verify all
reference metadata, and state the simulator's fidelity limits explicitly.
The most important limitations to discuss are the current low-fidelity
environment models, the chief-orbit approximation for distributed spacecraft
systems, and the distinction between internal MuJoCo forces and external
forces that feed back into orbit propagation.

\begin{thebibliography}{9}

\bibitem{gmat}
NASA Goddard Space Flight Center,
``General Mission Analysis Tool (GMAT).''
Available: \url{https://software.nasa.gov/software/GSC-17177-1}

\bibitem{orekit}
CS Group,
``Orekit: A low-level space dynamics library.''
Available: \url{https://www.orekit.org}

\bibitem{basilisk}
AVS Laboratory,
``Basilisk astrodynamics simulation framework.''
Available: \url{https://hanspeterschaub.info/basilisk/}

\bibitem{wertz1978spacecraft}
J.~R. Wertz, Ed.,
\emph{Spacecraft Attitude Determination and Control}.
D. Reidel Publishing Company, 1978.

\bibitem{wie2008space}
B.~Wie,
\emph{Space Vehicle Dynamics and Control}, 2nd~ed.
American Institute of Aeronautics and Astronautics, 2008.

\bibitem{todorov2012mujoco}
E.~Todorov, T.~Erez, and Y.~Tassa,
``MuJoCo: A physics engine for model-based control,''
in \emph{Proc. IEEE/RSJ International Conference on Intelligent Robots and Systems},
2012, pp.~5026--5033.

\bibitem{drake}
R.~Tedrake and the Drake Development Team,
``Drake: model-based design and verification for robotics.''
Available: \url{https://drake.mit.edu}

\bibitem{gazebo}
N.~Koenig and A.~Howard,
``Design and use paradigms for Gazebo, an open-source multi-robot simulator,''
in \emph{Proc. IEEE/RSJ International Conference on Intelligent Robots and Systems},
2004, pp.~2149--2154.

\end{thebibliography}


\end{document}
